\documentclass[10pt,letterpaper]{article}
\usepackage[utf8]{inputenc}
\usepackage[T1]{fontenc}
\usepackage[margin=0.75in]{geometry}
\usepackage{hyperref}
\usepackage{enumitem}
\usepackage{titlesec}
\usepackage{lmodern}

\pagestyle{empty}

\titleformat{\section}{\large\bfseries\uppercase}{}{0em}{}[\titlerule]
\titlespacing{\section}{0pt}{1ex}{1ex}

\newcommand{\CVItem}[1]{\item #1}

\begin{document}

\begin{center}
    {\LARGE \textbf{Jair Molina Arce}} \\ \vspace{0.5ex}
    Querétaro, Qro., México \\
    \href{mailto:ingjairmolina@gmail.com}{ingjairmolina@gmail.com} | 5652646108 \\
    \href{https://jairmolina.dev}{jairmolina.dev} | \href{https://github.com/smookymolina}{github.com/smookymolina} | \href{https://linkedin.com/in/jair-molina-arce-4909622b2}{linkedin.com/in/jair-molina-arce-4909622b2}
\end{center}

\vspace{1ex}

\begin{center}
    {\Large \textbf{Ingeniero Mecánico y de Sistemas Embebidos}}
\end{center}

\section*{Perfil Profesional}
Ingeniero Mecánico multidisciplinario con sólida experiencia transversal que abarca el diseño mecánico avanzado, el desarrollo de hardware (PCBs) y la programación de firmware de bajo nivel para sistemas embebidos. Mi trayectoria profesional incluye la caracterización detallada de bancos de prueba de propulsión, la ejecución de simulaciones térmicas y estructurales complejas (FEA), y la automatización integral de procesos físicos mediante microcontroladores (ESP32, STM32, arquitecturas ARM Cortex). Destaco en el sector por mi capacidad excepcional para actuar como puente tecnológico: integro fluidamente requerimientos mecánicos con soluciones electrónicas y de software, desarrollando productos desde la fase de diseño conceptual (CAD), pasando por el prototipado rápido y el diseño de circuito impreso, hasta llegar a la validación física del hardware final y su puesta en marcha.

\section*{Habilidades Técnicas Integrales}
\begin{itemize}[leftmargin=*, itemsep=0.2ex]
    \CVItem{\textbf{Diseño Mecánico y Simulación}: Dominio de SolidWorks para modelado 3D de piezas complejas y ensamblajes mayores; análisis por Elementos Finitos (FEA) empleando ANSYS para validación estructural bajo cargas estáticas/dinámicas y estrés térmico; modelado matemático en MATLAB/Simulink; metrología e interpretación de tolerancias geométricas (GD\&T).}
    \CVItem{\textbf{Desarrollo de Firmware y Electrónica}: Programación en C/C++ y MicroPython; arquitectura de software embebido (máquinas de estado, interrupciones, manejo de memoria); diseño de PCBs mixed-signal utilizando KiCad con enfoque en integridad de señal e inmunidad al ruido electromagnético; integración de periféricos (ADC, DAC, PWM) y buses (I2C, SPI, UART).}
    \CVItem{\textbf{Instrumentación y Control}: Diseño de sistemas DAQ (Adquisición de Datos); selección, calibración y acondicionamiento de sensores (temperatura, presión, celdas de carga); implementación y sintonización de lazos de control PID; manejo competente de equipo de laboratorio (osciloscopios, generadores de funciones, fuentes de poder).}
    \CVItem{\textbf{Software, Datos y Gestión}: Desarrollo de scripts en Python (Pandas, NumPy) para automatización de análisis de datos experimentales; desarrollo web backend (Flask); control de versiones (Git); uso de Docker para entornos aislados; redacción técnica y documentación de ingeniería.}
\end{itemize}

\section*{Experiencia Relevante}
\begin{itemize}[leftmargin=*, itemsep=0.5ex]
    \item \textbf{Docente Universitario e Investigador} | \textit{Instituciones Privadas y Públicas (IPN)} | 2024 -- Presente.
    \begin{itemize}[itemsep=0.1ex]
        \item Impartición de cátedras en Ingeniería Térmica, Modelización de Sistemas Aeroespaciales y Tecnologías Disruptivas.
        \item Desarrollo continuo de investigación en el área de control de sistemas dinámicos e instrumentación avanzada.
        \item Autoría de artículos técnicos y presentaciones en foros internacionales de ingeniería aeroespacial.
    \end{itemize}
    \item \textbf{Ingeniería y Desarrollo Integrado - Banco de Pruebas de Propulsión} | \textit{Instituto Politécnico Nacional} | 2022 -- 2024.
    \begin{itemize}[itemsep=0.1ex]
        \item Ejecuté el ciclo completo de diseño de una instalación experimental para motores cohete: desde el diseño CAD del soporte de empuje, hasta la selección de galgas extensiométricas y la programación del DAQ.
        \item Implementé rutinas de seguridad automatizadas mediante relés de estado sólido y electroválvulas controladas por STM32.
        \item Reduje drásticamente la incertidumbre de medición aplicando técnicas de filtrado digital directamente en el firmware.
        \item Consolidé una arquitectura de red local para la transmisión de datos críticos en tiempo real hacia la estación de control.
    \end{itemize}
\end{itemize}

\section*{Proyectos Destacados (Ciclo de Desarrollo Completo)}
\begin{itemize}[leftmargin=*, itemsep=0.5ex]
    \item \textbf{Jaguar MX - Sistema de Extracción de Aire Caliente para Gabinete de Telecomunicaciones} (2026).
    \begin{itemize}[itemsep=0.1ex]
        \item Desarrollo end-to-end de un sistema de gestión térmica industrial. Inicié con el cálculo de cargas térmicas y diseño CAD para el montaje de ventiladores y ductos. Posteriormente, diseñé el hardware (esquemático y PCB en KiCad) aislando la lógica de baja tensión del puente H de potencia. Finalmente, programé el firmware en C++ asegurando un funcionamiento ininterrumpido mediante watchdogs de hardware y máquinas de estados robustas contra fallos de sensores.
    \end{itemize}
    \item \textbf{Impresora 3D DIY - Diseño Mecatrónico desde Cero}.
    \begin{itemize}[itemsep=0.1ex]
        \item Proyecto integral que combinó el diseño CAD del pórtico, análisis de deflexiones, selección de perfiles estructurales, cableado eléctrico del cuadro de control, flasheo y configuración de firmware Marlin, y calibración final de pasos por milímetro e interpolación térmica del hotend.
    \end{itemize}
    \item \textbf{Sistema de Telemetría Inalámbrica y Control Adaptativo PID}.
    \begin{itemize}[itemsep=0.1ex]
        \item Desarrollo de un nodo sensor portátil con comunicación WiFi/MQTT, operando con baterías Li-Po. Implementación de algoritmos de control adaptativo para compensar variaciones de inercia térmica en tiempo real, mejorando la respuesta del sistema frente a perturbaciones externas.
    \end{itemize}
    \item \textbf{Pipeline de Automatización de Análisis de Datos de Ingeniería}.
    \begin{itemize}[itemsep=0.1ex]
        \item Software propietario desarrollado en Python que procesa archivos de registro (logs) generados durante pruebas destructivas, detecta anomalías, aplica filtros de media móvil y genera automáticamente reportes ejecutivos en formato PDF y hojas de cálculo con gráficos interactivos.
    \end{itemize}
\end{itemize}

\section*{Educación}
\begin{itemize}[leftmargin=*, itemsep=0.2ex]
    \CVItem{\textbf{Maestría en Tecnologías Avanzadas} (En curso) - IPN.}
    \CVItem{\textbf{Ingeniería Mecánica} (2017 -- 2023) - IPN.}
    \CVItem{\textbf{Bachillerato Técnico en Informática} (2014 -- 2017) - Colegio de Bachilleres.}
\end{itemize}

\section*{Idiomas y Reconocimientos}
\begin{itemize}[leftmargin=*, itemsep=0.2ex]
    \CVItem{\textbf{Idiomas}: Español (Nativo), Inglés (Avanzado, B2/C1 técnico).}
    \CVItem{\textbf{Publicaciones}: Coautor presentado en el \textit{75th IAC 2024} (Italia). Artículo en \textit{Revista Hacia el Espacio}.}
\end{itemize}

\end{document}
